\section{Gradient preprocessing}\label{sec:pre}

One potential challenge in training optimizers is that different input coordinates (i.e. the gradients w.r.t. different
optimizee parameters) can have very different magnitudes.
This is indeed the case e.g. when the optimizee is a neural network and different parameters
correspond to weights in different layers.
This can make training an optimizer difficult, because neural networks
naturally disregard small variations in input signals and concentrate on bigger input values.

To this aim we propose to preprocess the optimizer's inputs.
One solution would be to give the optimizer $\left(\log(|\nabla|),\,\operatorname{sgn}(\nabla)\right)$ as an input, where
$\nabla$ is the gradient in the current timestep.
This has a problem that $\log(|\nabla|)$ diverges for $\nabla \rightarrow 0$.
Therefore, we use the following preprocessing formula 
\begin{align*}
    \nabla^k &\rightarrow
\begin{cases}
    \left(\frac{\log(|\nabla|)}{p}\,, \operatorname{sgn}(\nabla)\right)& \text{if } |\nabla| \geq e^{-p}\\
    (-1, e^p \nabla)             & \text{otherwise}
\end{cases}
\end{align*}
where $p>0$ is a parameter controlling how small gradients are disregarded (we use $p=10$ in all our experiments).

We noticed that just rescaling all inputs by an appropriate constant instead also works fine, but
the proposed preprocessing seems to be more robust and gives slightly better results on some problems.



%Different optimizees $f$ have very different gradient magnitudes, which can make training an optimizer very difficult.
%One possible way to mitigate this problem is Batch Normalization \todo{cite}.
%It however has a disadvantage that it assumes that the distribution of the gradients is stationary.

%Therefore, we propose 2 different solutions for this problem.
%The simplest one is to simply scale all the gradients by some constant (the same for all timesteps), which makes them
%have a distribution with a standard deviation approximately $1$.

%The way neural networks are trained naturally causes them to disregard
%very small variations in input signals.

%We have found that it can be useful to transform the input to the
%optimizer network into a representation that takes this into account. For example we have considered the transformation,


